\section{Workloads}
\label{sec:workloads}

Every workload implements one interface: \code{prepare} (untimed, before
each iteration), \code{run} (the timed work) and \code{checksum} (untimed).
Buffers are allocated in the constructor: WASM's default allocator takes a
global lock, so an allocation inside the timed region would add a
synchronization cost that native code does not have.

\paragraph{Inputs.} Images are 8192$\times$8192 single-channel \code{u8}
pixels (67 megapixels, 64~MiB), generated outside the timed region by a
SplitMix64 generator implemented in the crate and seeded with a fixed
constant, so the pixels are bit-identical in every environment. The
\emph{noise} input takes uniform values 0 to 255. The \emph{few-colours}
input takes $k$ values evenly spaced over 0 to 255, each pixel picking one
uniformly. Every run records the seed, $k$ and a fingerprint of the input.

\paragraph{Gaussian filter: no synchronization.} A $3\times3$ kernel
$[1\,2\,1; 2\,4\,2; 1\,2\,1]/16$ in integer arithmetic with rounding, edges
clamped. Rows are processed in parallel (\code{par\_chunks\_mut} over the
output rows): each task writes its own row and reads the shared, immutable
source, so threads never coordinate, and any loss of speedup comes from the
environment. Integer arithmetic gives the same result on every target and
thread count, where floating point could differ with SIMD. The interior
loop works on 16-bit values (the largest sum, $16\times255$, fits) and is
the function whose vectorization is checked. The filter reads one 64~MiB
image and writes another.

\paragraph{Histogram: two variants with the same result.} The histogram
counts the pixels of each value into 256 bins.
\begin{itemize}[leftmargin=1.5em]
\item \emph{Atomic}: all threads increment 256 shared \code{AtomicU64}
counters (\code{fetch\_add(1)}). Each counter is padded to its own 128-byte
block, so that two different counters never share a cache line: what is
measured is true sharing, not false sharing. 128 and not 64 bytes because
the adjacent-line prefetcher fetches 64-byte lines in pairs. The counters
are reset in \code{prepare}, outside the timed region.
\item \emph{Local}, the lock-free control: each task counts into its own
plain table and the tables are merged at the end (\code{fold} and
\code{reduce}). No write is shared.
\end{itemize}
Both split the pixels into 64~KiB chunks and build their tasks through the
same code, both index fixed-size tables (no bounds check per pixel), and
both time their final step (the reduce, or the read-back of the shared
counters). They differ in how the counts are combined, which is the factor
of RQ2. The memory ordering asked for is \code{Relaxed}; it makes no
difference on these targets, since x86 locked instructions are full
barriers and WASM atomics are always sequentially consistent.

\paragraph{Contention is set by $k$.} The histograms run on the few-colours
input with $k \in \{1, 4, 16, 256\}$: the updates are spread over $k$
counters, from every pixel on one counter to a uniform spread over all 256.
The first design called the two ends ``high'' and ``low'' contention; a
review before any data replaced the names, because with 256 counters and
eight threads most increments still find their line modified by another
core. What $k$ changes is how the contention is spread. $k$ also changes
the control: with few values the local histogram increments the same
address again and again, a chain in which each increment waits for the
previous store (store-to-load forwarding), with no synchronization
involved. The RQ2 metric is built so that this effect of the input cancels
(Section~\ref{sec:plan}).

\paragraph{Padding, checked.} A pilot compared 64-byte with 128-byte padding
natively (10 runs per cell): every one-thread cell within 0.6\%, and at
eight threads ratios of 0.99 to 1.02 with intervals that reach 7\% either
way at $k=1$; no difference was detected and 128 was kept as the safer
choice. A \code{perf c2c} run in Chrome at $k=16$ showed 16 contended lines
0x800 bytes apart (16 bins of 128 bytes), each accessed at one offset: the
alignment holds inside WASM linear memory. The 64-byte variant was later
removed from the code.

\paragraph{Checksums.} Every workload hashes its output (FNV-1a, which
depends on the position of every byte) outside the timed region. The
checksum of a run must equal that of every other run of the same workload
and input, whatever the environment and the thread count; a conflict stops
the analysis. Golden checksums are pinned in the tests, and the continuous
integration runs every workload in both browsers and compares checksums and
input fingerprints with the native run. No campaign of v3 has a conflict.

\paragraph{Sequential references.} \code{filter-seq} and \code{hist-seq} are
the same kernels called from a plain loop with no Rayon, at one thread
only. They answer whether a pool of one thread is a fair baseline for the
speedup (Section~\ref{sec:sequential}).

\paragraph{Variants of the exploratory campaigns.} \emph{Private} atomics:
every thread increments its own padded table of atomic counters, so the
atomic instruction is paid and no line is shared; with local and shared it
forms a ladder that separates the cost of the instruction from the cost of
contention. \emph{Packed} atomics: the 256 shared counters are contiguous,
eight per 64-byte line, the false-sharing level next to the padded one.
\emph{Sorts}: Rayon's stable parallel sort, a merge sort that allocates its
buffer and so exercises WASM's allocator, and the unstable in-place sort as
its control, on the pixels read as 32-bit keys.
